\scanpage{326}\mspage{50}
% archive sheet number: 325

$W_0=W$, on trouve donc l'implication $A\Longrightarrow D$, qui est \deleted{$A\unsure{=|\Longrightarrow}$} également contenue dans les autres implications du diagramme, comme comprise dans la chaîne d'implications
\[
 A\Longrightarrow B\Longrightarrow C\Longrightarrow D\Longrightarrow D_{W_0}
 \overset{\bullet}{\Longrightarrow}F.
\]

Il reste, dans ces commentaires généraux, à dire un mot de cette dernière implication
\[
 D_{W_0}\overset{\bullet}{\Longrightarrow}F,
\]
i.e. du fait que la plus faible de toutes les \deleted{\unsure{les|de}} conditions \deleted{\unsure{à|de} \illegible{1}} $A$ à $D_{W_0}$ envisagées ici \deleted{(à la seule exception de $B_{W_0}$ et $C_{W_0}$, qui \unsure{même})} implique déjà que $\mathrm{Cat}$ a \unsure{``}assez de $W$-fibrations, i.e est une catégorie à \unsure{cofibrations} pour $(W,\deleted{\unsure{C}}\mathrm{Fib}_{W})$. C'est clair quand on prend $D_{W_0}$ sous la forme $A(\mathrm{Parf}_{W_0})$ i.e
\[
 \unsure{\mathrm{Fi}|\mathrm{F}}\hspace{3em}\mathrm{Parf}_{W_0}\subset\unsure{\ldots W|\mathrm{Fib}_{W}}.
\]
Comme \unsure{tte|toute} flèche $f$ de $\mathrm{Cat}$ se factorise en $f=pi$, avec $\deleted{p\in}i\in W_0\subset W$ et $f\in\mathrm{Parf}_{W_0}\subset\mathrm{Fib}_{W}$, on trouve la factorisation avec $i\in W$, $f\in\mathrm{Fib}_{W}$.
